\section{Como o sistema funciona}
O projeto guarda os mesmos dados em três réplicas, cada uma com seu próprio
arquivo JSON. O cliente envia os pedidos a um coordenador, que decide onde
ler e gravar conforme o modo escolhido: strong, eventual ou Read Your Writes
(RYW, que significa ``ler as próprias escritas''). A comunicação usa HTTP.

O coordenador foi implementado em \texttt{coordenador.py}, e cada réplica
executa \texttt{replica.py}. Na configuração manual, o coordenador atende na
porta 5000 e as réplicas nas portas 5001, 5002 e 5003. Uma escrita envia a
chave e o valor por \texttt{POST /write}; uma leitura usa
\texttt{GET /read?chave=...}. O cliente acessa o coordenador, que encaminha
esses pedidos às réplicas.

Cada registro contém o valor e um número de versão. Para ler, a réplica abre
o JSON e procura a chave. Para escrever, lê o arquivo inteiro, altera o
registro e grava novamente o conteúdo. Primeiro, salva em um arquivo
temporário, usa \texttt{flush} e \texttt{fsync} para enviar os dados ao
armazenamento e depois substitui o arquivo antigo com \texttt{os.replace}.
Isso evita deixar um JSON gravado pela metade. Um bloqueio organiza o acesso
quando chegam várias operações ao mesmo tempo.

\section{Como o eventual foi implementado}
O eventual permite que as réplicas tenham valores diferentes por um tempo,
enquanto as atualizações são enviadas. O caminho de uma escrita é:
\begin{enumerate}
\item O coordenador escolhe uma réplica por rodízio: 1, 2, 3, 1, e assim por diante.
\item Cria uma versão para a escrita e envia a chave, o valor e a versão à réplica.
\item Se ela confirmar a gravação, coloca uma tarefa na fila para atualizar as outras duas.
\item Responde sucesso ao cliente, sem esperar essas duas atualizações.
\end{enumerate}

As leituras também alternam entre réplicas, usando o mesmo contador de
rodízio das escritas. Elas consultam somente a cópia escolhida, sem comparar
com as demais. Por isso, uma leitura logo após uma escrita pode devolver
o valor antigo ou até 404, caso a chave ainda não tenha chegado àquela réplica.

Por exemplo, se todas as cópias guardam 50 e uma escrita muda a réplica 1
para 80, o cliente já pode receber sucesso. Uma leitura da réplica 2 ainda
pode devolver 50 até ela ser atualizada. Depois da propagação, as três terão
80, desde que outra escrita mais nova não tenha substituído esse valor.

\subsection{Fila e tempo de sincronização}

A sincronização usa uma fila de prioridade em memória, implementada com
\texttt{heapq}. Cada tarefa guarda a chave, o valor, a versão e os destinos.
O horário em que a tarefa pode começar define sua posição na fila. Uma única
thread de replicação espera esse horário e envia as atualizações por HTTP,
um destino por vez. As réplicas não trocam dados diretamente entre si.

O atraso é definido por \texttt{--replication-delay}. Seu valor padrão,
também usado nos testes, é de 1 segundo,
contado a partir do agendamento após a tentativa de escrita inicial.
Esse é o tempo até o envio poder começar, e não um prazo para as três cópias
ficarem iguais: a fila, os bloqueios e as gravações podem aumentar a espera.

O agendamento usa \texttt{time.monotonic()}, um relógio usado para medir
intervalos. O programa não espera mais um segundo inteiro entre todas as
tarefas: se várias já estiverem prontas, pode atendê-las em sequência.
Com atraso zero, o envio começa assim que houver capacidade, mas continua
sendo feito em segundo plano.

Mesmo no eventual, as escritas iniciais passam por um bloqueio global.
A thread de replicação não usa esse bloqueio nesse modo, mas disputa os
bloqueios de cada réplica com as outras operações. Assim, aumentar a carga
pode fazer a fila crescer e atrasar a sincronização.

\subsection{Versões e tratamento de falhas}

Se um envio falhar, ele é agendado novamente, neste caso com atraso de
1 segundo. Somente os destinos que falharam entram nessa nova tentativa.
A regra geral é esperar o maior valor entre 0,1 segundo e o atraso configurado.
O timeout de comunicação com cada réplica é de 2 segundos. Esse limite não
inclui toda a espera na fila e nos bloqueios, portanto uma operação completa
pode demorar mais.

A versão de uma escrita é o maior valor entre o horário atual em nanossegundos
e a última versão criada mais um. Isso faz as versões crescerem durante a
execução do coordenador. A réplica recusa uma versão menor que a já salva.
Se receber novamente a mesma versão com o mesmo valor, confirma sem regravar;
se a versão for igual, mas o valor for diferente, retorna conflito.
Essa proteção evita que um envio atrasado substitua um dado mais novo.

As tentativas continuam enquanto a escrita ainda for a mais nova da chave
e o coordenador estiver ativo. Quando chega uma escrita mais nova para a
mesma chave, as tarefas antigas deixam de ser propagadas, e a nova passa
a ser responsável pela atualização.

Se a própria escrita inicial falhar, o cliente recebe 503 e o coordenador
agenda o envio para todas as réplicas, inclusive a origem. Isso acontece
porque um timeout pode ocorrer depois de a réplica já ter gravado o valor.
Portanto, receber erro não significa que a escrita foi cancelada: ela pode
ser concluída pelas novas tentativas.

A fila fica em memória e pode ser perdida se o coordenador for encerrado.
O programa não reconstrói essa fila ao reiniciar e não faz uma comparação
periódica de todos os arquivos para corrigir qualquer diferença.

\section{Como o RYW foi implementado}
O RYW usa a mesma sincronização do eventual, mas mantém cada cliente em uma
réplica fixa. O coordenador usa o identificador enviado no cabeçalho
\texttt{X-Client-ID} para escolher essa réplica por meio de um hash SHA-256.
O mesmo identificador e a mesma lista ordenada de réplicas levam ao mesmo destino.

O código calcula SHA-256 do identificador, transforma os primeiros oito bytes
em um número inteiro e usa o resto da divisão pela quantidade de réplicas.
Com três réplicas, o resultado escolhe uma das posições 0, 1 ou 2 da lista.
Clientes diferentes podem cair na mesma réplica; a distribuição não precisa
ficar perfeitamente equilibrada.

Depois dessa escolha, a escrita segue o mesmo processo do eventual: confirma
na réplica escolhida, agenda as outras e responde. O RYW não usa uma fila
mais rápida nem um tempo de sincronização diferente. A mudança está em
manter a escrita e a leitura do cliente no mesmo destino.

Como a escrita e a leitura do cliente vão para a mesma cópia, depois de uma
escrita confirmada ele consegue ler o que escreveu, ou um valor mais novo.
Outros clientes ainda podem encontrar cópias atrasadas. O identificador é
obrigatório e deve ser mantido entre os pedidos. Se a réplica fixa estiver
indisponível, o sistema retorna erro em vez de consultar uma cópia possivelmente
desatualizada.

Por exemplo, se o cliente A estiver associado à réplica 2, ele grava 80 nessa
cópia e sua próxima leitura também consulta a réplica 2. Já um cliente B,
associado à réplica 3, pode continuar lendo 50 até a atualização chegar.
Isso mostra que o RYW protege a sequência de operações de cada cliente,
mas não faz todos os clientes enxergarem imediatamente a mesma informação.

O código não guarda uma tabela com a última versão de cada cliente. A
garantia vem da réplica fixa, da confirmação da escrita nela e da recusa
de versões antigas. Sem \texttt{X-Client-ID}, o coordenador retorna 400.
Trocar o identificador ou a ordem da lista de réplicas pode mudar o destino
e comprometer essa garantia. No YCSB, cada thread usa um identificador UUID
estável nos seus pedidos, permitindo testar esse modo.

\section{Como os testes foram feitos}
Os resultados salvos foram obtidos com YCSB, usando 1.000 registros,
1.000 operações por execução e oito threads. Foram dez rodadas para cada
combinação de modo e carga: leitura, escrita e mista, com 50\% de probabilidade
para cada operação. Isso totaliza 90 execuções e 90.000 operações.
Cada registro tinha um campo de 100 caracteres, com escolha uniforme das chaves.

Todos os processos rodaram no mesmo computador, com Windows 11, Python 3.13.7
e Java 21.0.7. Cada caso começou com uma cópia da mesma base, já sincronizada,
e novos processos. A carga inicial ficou fora da medição. Não houve uma etapa
separada de aquecimento antes dos testes.

O binding, que liga o YCSB à API do coordenador, transforma INSERT e UPDATE
em pedidos de escrita e READ em pedidos de leitura. A carga medida de escrita
usa UPDATE, ou seja, atualizações de registros existentes. No teste misto,
as operações são sorteadas: 50\% é a probabilidade, não uma exigência de
exatamente 500 leituras e 500 escritas em cada rodada.

Os casos foram executados um por vez, alternando a ordem dos modos entre
rodadas. As médias dão o mesmo peso a cada rodada. O throughput indica
operações por segundo e inclui as que retornaram erro. Já as latências
apresentadas são os tempos das operações bem-sucedidas. O desvio-padrão
(DP) mostra quanto os resultados variaram de uma rodada para outra.

\section{Resultados e comparações}

\begin{table}[htbp]
\centering\small
\setlength{\tabcolsep}{4pt}
\caption{Operações por segundo: média $\pm$ DP; erros nas 30.000 operações de cada modo.}
\begin{tabular}{llllr}
\toprule
Modo & Leitura & Escrita & Mista & Erros\\\midrule
strong & 41,79 $\pm$ 7,20 & 7,87 $\pm$ 1,46 & 12,46 $\pm$ 1,46 & 30\\
eventual & 498,35 $\pm$ 103,50 & 21,22 $\pm$ 3,89 & 40,50 $\pm$ 3,97 & 56\\
ryw & 383,85 $\pm$ 101,94 & 19,96 $\pm$ 3,62 & 43,58 $\pm$ 8,48 & 36\\
\bottomrule
\end{tabular}
\end{table}

Quanto maior a quantidade de operações por segundo, maior a taxa de atendimento.
O eventual teve a maior média em leitura e escrita, enquanto o RYW teve a
maior média na carga mista.

\textbf{O strong demorou mais porque espera mais confirmações antes de
responder e limita mais as operações simultâneas.} Uma escrita nele só
recebe sucesso após a confirmação das três réplicas. Uma leitura consulta
as três e compara seus valores. Os pedidos às réplicas são feitos um após
o outro, e um bloqueio global faz as leituras e escritas do coordenador
esperarem umas pelas outras. Eventual e RYW aguardam uma réplica na escrita
e consultam uma na leitura, deixando a propagação restante em segundo plano.

O JSON também pesa nessa diferença: a escrita regrava o arquivo inteiro,
e o strong espera esse trabalho nas três cópias. Se uma escrita falhar
parcialmente, ele retorna erro, tenta sincronizar novamente e bloqueia a
leitura daquela chave até a recuperação. Leituras com valores diferentes
entre as réplicas retornam conflito. Essa proteção depende de um único
coordenador e do acesso dos clientes por ele, com a base inicialmente coerente.

\subsection{Leitura}
O eventual atingiu 498,35 ops/s, contra 383,85 no RYW e 41,79 no strong.
Em relação ao strong, isso representa cerca de 11,93 vezes a taxa no eventual
e 9,19 vezes no RYW. O tempo médio de resposta foi de 196,90 ms no strong,
16,31 ms no eventual e 18,27 ms no RYW.

Consultar uma única cópia ajuda a explicar o ganho dos modos assíncronos.
O eventual ficou cerca de 29,83\% acima do RYW na média. O rodízio do eventual
e o destino fixo por cliente do RYW podem distribuir a carga de formas
diferentes, mas os testes não mediram esse efeito separadamente. Como os
identificadores são distribuídos por hash, as oito threads do RYW não precisam
ficar igualmente divididas entre as três réplicas.

\subsection{Escrita}
O eventual atingiu 21,22 ops/s e o RYW 19,96 ops/s, contra 7,87 ops/s no strong.
São aproximadamente 2,70 e 2,54 vezes a taxa do strong. Os tempos médios
de resposta foram 1.052,94 ms no strong, 385,44 ms no eventual e 409,94 ms no RYW.

O ganho ocorre porque o cliente recebe a resposta após uma confirmação,
enquanto as outras cópias ainda serão atualizadas. Isso melhora o tempo de
resposta, mas não elimina o custo de manter as três réplicas: parte dele
aparece na espera pela sincronização após o teste.

\subsection{Carga mista}
O RYW atingiu 43,58 ops/s, o eventual 40,50 ops/s e o strong 12,46 ops/s.
As taxas do RYW e do eventual foram cerca de 3,50 e 3,25 vezes a do strong.
O RYW ficou 7,59\% acima do eventual na média, mas também variou mais:
seu DP foi de 8,48 ops/s, contra 3,97 ops/s no eventual. Essa diferença
não basta para afirmar que ele sempre será melhor.

O tempo médio de leitura foi de 623,18 ms no strong, 52,17 ms no eventual
e 45,09 ms no RYW. Para a escrita, foi de 674,19 ms, 325,39 ms e 316,35 ms,
respectivamente. O aumento da latência de leitura do strong em relação ao
teste só de leitura é compatível com a espera por escritas no bloqueio global.
Essa é uma explicação baseada no código; o tempo gasto no bloqueio não foi
medido isoladamente.

\subsection{Tempo para as réplicas ficarem iguais}
As três réplicas ficaram iguais ao final de todas as 90 execuções.
Depois do teste de escrita, a espera média adicional foi de 0,13 s no strong,
26,72 s no eventual e 23,22 s no RYW. Na carga mista, foi de 0,02 s,
17,71 s e 15,57 s, respectivamente. Isso mostra por que o atraso configurado
de 1 segundo não deve ser confundido com o tempo total de sincronização.

Essa espera foi medida após o YCSB terminar, comparando os valores e versões
dos arquivos, com novas verificações a cada 0,25 s quando havia diferenças.
Ela não entra nas taxas de operações por segundo e não mede o atraso de cada
escrita. Uma parte do ganho dos modos assíncronos vem de responder ao cliente
antes de terminar todo o trabalho de replicação.

No teste somente de escrita, a duração média da fase medida foi de 132,38 s
no strong, 48,34 s no eventual e 51,43 s no RYW. Somando a espera posterior,
os tempos ficam em aproximadamente 132,51 s, 75,06 s e 74,65 s. Essa conta
ajuda a comparar o trabalho até a igualdade observada das cópias, mas não
inclui toda a preparação dos processos e não substitui o throughput do YCSB.

\subsection{Erros e limites da comparação}
O total de erros foi de 30 no strong, 56 no eventual e 36 no RYW, em
30.000 operações por modo. As taxas do YCSB incluem operações que falharam.
Esses erros, sozinhos, não mostram quantas leituras devolveram valores antigos.
Da mesma forma, encontrar as cópias iguais ao final não prova que todas as
leituras durante o teste respeitaram as garantias de consistência.

Na carga de leitura, houve 10 erros no strong, 45 no eventual e 23 no RYW.
Na escrita, foram 20, 11 e 12; na carga mista, zero, zero e um, respectivamente.
A base da carga somente de leitura já estava sincronizada e não recebia
escritas, portanto seus erros não devem ser tratados como medida de atraso
de replicação. O binding distingue respostas como NOT\_FOUND e
SERVICE\_UNAVAILABLE, mas falhas de comunicação também podem virar ERROR.
Os totais não identificam sozinhos a causa de cada falha.

Os testes foram locais, em sessões diferentes, sem controle completo da carga
do computador e sem simulação planejada de falhas de rede. O atraso de
1 segundo foi aplicado ao replicador, não à rede inteira. Os resultados
descrevem esse ambiente e esse armazenamento em JSON; não devem ser tratados
como desempenho geral de qualquer sistema com esses modelos de consistência.

Os resultados são de testes já existentes. Os identificadores do conteúdo
dos arquivos de código (hashes) atuais diferem dos registrados nesses testes.
Por isso, as explicações descrevem o código disponível, mas não foi possível
confirmar que ele é exatamente a mesma versão usada nas medições.

\section{Conclusão}
O strong teve menos operações por segundo e respostas mais lentas porque
espera as três cópias e organiza as operações uma de cada vez no coordenador.
O eventual e o RYW respondem após uma confirmação e sincronizam o restante
em segundo plano, mas podem deixar trabalho pendente depois do teste.
O eventual pode devolver um valor antigo após uma escrita; o RYW evita isso
para o mesmo cliente ao mantê-lo na réplica onde sua escrita foi confirmada.
A escolha depende de quanto a aplicação precisa que os dados estejam
atualizados no momento da leitura.
